\par\medskip\noindent\textbf{Cómo leer la salida.}
{\small\setlength{\tabcolsep}{4pt}\renewcommand{\arraystretch}{1.15}
\begin{xltabular}{\linewidth}{@{}>{\raggedright\arraybackslash}p{0.24\linewidth}>{\raggedleft\arraybackslash}p{0.15\linewidth}>{\raggedright\arraybackslash}X@{}}
\toprule Indicador & Valor & Interpretación \\ \midrule \endfirsthead
\toprule Indicador & Valor & Interpretación \\ \midrule \endhead
\bottomrule \endlastfoot
\texttt{T} & 929 meses & Muestra mensual 1946--2023. \\
\texttt{b} & 0,00611 ($t=2{,}35$) & Coeficiente de predicción de $\ln(D/P)$ sobre el rendimiento del mes siguiente; parece significativo. \\
\texttt{phi} & 0,99759 & Persistencia mensual del predictor: casi una raíz unitaria. \\
\texttt{corr(u, v)} & $-0{,}983$ & Correlación entre el error de la regresión de rendimientos y la innovación del predictor: casi perfecta y negativa (si sube el precio, baja $D/P$). \\
$\sigma_{uv}/\sigma_v^2$ & $-0{,}9541$ & Factor que traslada el sesgo de $\hat\phi$ al de $\hat b$. \\
Sesgo de Kendall & $-0{,}00430$ & Sesgo hacia abajo de $\hat\phi$ en muestra finita, $\approx-(1+3\phi)/T$. \\
Sesgo de Stambaugh & 0,00410 & Sesgo hacia arriba de $\hat b$ que resulta de los dos anteriores. \\
\texttt{b corregido} & 0,00201 & Coeficiente sin sesgo: un tercio del estimado. \\
Proporción atribuible al sesgo & 67,1\,\% & Dos tercios de la ``predictibilidad'' estimada son sesgo de muestra finita. \\
\end{xltabular}}
